\chapter{The story in words}\label{ch:story}

\begin{chapterplan}
The question of this chapter: what does the paper claim, and why, before any mathematics?
After this chapter you can do five things.
\begin{itemize}
\item Say why the usual answer to the paper's question is ``no'', and what the paper adds.
\item Name each player and say what each one knows, and when.
\item Explain the two returns to information that a stronger incumbent changes.
\item Tell the chain of the main result, and say where the chain breaks against a very strong incumbent.
\item Say which chapter of this guide proves each claim.
\end{itemize}
\end{chapterplan}

This chapter tells the whole paper in words. It has no derivation. Every later chapter proves one part of the story, so each derivation has a place that you already know. A symbol appears only after the text names the object that it stands for. Later chapters define every symbol again and state each result exactly.

\section{The question}\label{sec:story-question}

A company is for sale. One buyer has already done the work that a serious offer needs. This buyer is the incumbent. A second buyer, the challenger, has not yet decided whether to do that work. The paper asks one question: can a stronger incumbent bring the challenger into the contest?

The usual answer is no. A challenger must prepare before it can make an offer that the seller can accept. It must check the target's contracts, debts and technology. It must also arrange money and approvals, and it must set aside staff time. These costs are sunk once the challenger pays them. So the challenger prepares only if the profit that it expects from the contest covers the cost.

A stronger incumbent lowers that profit in two ways. The challenger wins less often. When the challenger wins, it pays the incumbent's bid, and a stronger incumbent bids more. With less profit, fewer challengers find preparation worth its cost. The literature calls this force deterrence (Fishman 1988; Hirshleifer and Png 1989). If the challenger's information stays the same, a stronger incumbent can only lower entry.

The paper studies a force in the opposite direction. The target is a listed company, so its stock trades while the contest is open. The challenger can see the stock price before it decides whether to prepare. If the price carries news about how well the challenger fits the target, the challenger can act on the news. The paper shows that a stronger incumbent can make the price carry such news. The information effect can then outweigh deterrence.

\begin{keyidea}
A stronger incumbent has a direct effect and an indirect effect on the challenger. The direct effect lowers what preparation earns. The indirect effect makes the stock price more informative about the challenger. The challenger reads the price before it decides.
\end{keyidea}

Two terms recur in this guide. Entry is the probability that the challenger prepares. High-class ownership is the probability that a challenger from the high class acquires the target. The next section explains what a class is.

{\small\color{notegrey}Source: Section 1 of the paper.}

\section{The set-up}\label{sec:story-setting}

The model describes a listed target whose sale is public while its stock trades. A further bidder can still enter. The paper calls this period the decision interval. A disclosed approach, an announced strategic review or an open contest can start it. Section 2.1 of the paper describes the interval, and Chapter~\ref{ch:empirical} shows how to date it in data.

Six parties take part. The list says what each one does and knows.
\begin{itemize}
\item \emph{The target and its seller.} Before any trade, the seller commits to a cash second-price auction with a reserve price. Each prepared bidder bids its own value. The highest bid wins, and the winner pays the larger of the reserve and the highest other bid. The target receives this payment, which the paper calls the target proceeds.
\item \emph{The incumbent.} It has prepared. It knows its own value for the target before it bids. The public knows only the distribution of this value. A stronger incumbent is one whose distribution moves up, so that high values become more likely.
\item \emph{The challenger.} It has not prepared. Preparation costs money, and it shows the challenger its value for the target. The cost is zero with probability \(\rho\) and a high cost \(c_H\) otherwise. The challenger learns its cost after it sees the price. It belongs to class \(H\) or class \(L\), and the class sets the distribution of its value. Before it prepares, it knows neither its class nor its value.
\item \emph{The informed investor.} It is an outside specialist, such as a sector fund. It holds no stake in the target and cannot buy the target. It knows the challenger's class, but not its value. It earns only the profit from its trade.
\item \emph{The liquidity trader.} This trader acts for reasons that have nothing to do with the sale. It sells one unit, buys one unit, or does nothing.
\item \emph{The market makers.} They are competitive and risk neutral. They see the order and nothing else. They set the price equal to the expected payoff of a share, and they absorb the order at that price.
\end{itemize}

The two classes are ranked. At every level of value, a class-\(H\) challenger is at least as likely as a class-\(L\) challenger to value the target above that level. Every value also meets the reserve. The paper calls these two restrictions the class ranking. The prior is the probability, before any trade, that the challenger is of class \(H\).

Trade takes place in one order slot, a single public order opportunity before the challenger decides. With some probability the informed investor fills the slot, and otherwise a liquidity trader fills it. The order is one unit bought, one unit sold, or no order. An order to buy or to sell costs the investor a trading cost. Nobody sees who filled the slot. The slot follows Dow and Gorton (1997), where an informed order and a liquidity order are alternative arrivals in one opportunity.

The events happen in this order.
\begin{enumerate}[nosep]
\item The seller announces the sale rule.
\item The investor learns the challenger's class.
\item The slot is filled, and the order arrives.
\item The market makers see the order and set the price.
\item The challenger sees the price and learns its cost. It decides whether to prepare.
\item The prepared bidders learn their values and bid. The sale takes place.
\end{enumerate}

\begin{watchout}
The challenger sees the price and its own cost, but not the order. The investor sees the class, but not the value. The price can carry news only through what the market makers infer from the order.
\end{watchout}

The paper studies equilibrium. Each player's choice is a best reply to the choices of the others, and each price is the market makers' expected payoff of a share. Chapter~\ref{ch:model} states the six conditions. One term from there recurs. The entry floor \(\rho\) is the least entry that any price can produce, because a zero-cost challenger prepares after every price. The posterior is the market makers' belief that the class is \(H\), after they see the order.

{\small\color{notegrey}Source: Section 2 of the paper.}

\section{The two returns to information}\label{sec:story-returns}

Information has a value to two parties. The challenger uses it to decide whether to prepare. The investor uses it to trade. A stronger incumbent changes the two values in opposite directions. Proposition~1 of the paper states how.

The challenger holds a belief that its class is \(H\). Its gross acquisition profit is what it expects to keep from a win, before it pays its cost. If class \(H\) expects more than class \(L\), gross profit rises with the belief. The high cost fixes one belief at which gross profit equals this cost, and the paper calls this belief the preparation threshold. A high-cost challenger prepares when its belief reaches the threshold. It stays out when its belief is lower.

The target's proceeds depend on who prepares. A challenger of class \(H\) raises the price more than a challenger of class \(L\), because its value is likelier to be high. The difference between the two classes in the target's expected proceeds is the target-payoff spread, written \(\DT\). The spread is what the investor's knowledge is worth. The investor knows the class, and a wider spread means that the class matters more for the value of a share.

Proposition~1 makes two statements about a stronger incumbent.
\begin{itemize}
\item[(a)] It never raises the challenger's gross profit at any belief. This holds for every pair of classes.
\item[(b)] Under the class ranking, it never narrows the target-payoff spread.
\end{itemize}

Both statements come from one picture, the picture of slices. Cut the gap between the two classes into thin slices, one for each level of value above the reserve. The height of a slice is the gap between the two classes in the chance that the challenger's value lies above that level. The class ranking makes every slice non-negative.

Fix one incumbent value. In a second-price auction, the part of the challenger's value below the incumbent's value goes to the target. The part above stays with the challenger. So the slices to the left of the incumbent's value go to the target, and the slices to the right stay with the challenger. Figure~\ref{fig:story-slices} shows this.

\begin{figure}[t]
\centering
\begin{tikzpicture}[x=1cm,y=1cm,font=\footnotesize]
\foreach \c/\xs/\lab in {2.0/0/{(a) A lower incumbent value}, 3.7/7.6/{(b) A higher incumbent value}} {
  \begin{scope}[xshift=\xs cm]
    \begin{scope}
      \clip (-0.1,-0.1) rectangle (\c,3.2);
      \fill[accent!40] (0,0) -- plot[domain=0:6,samples=40] (\x,{3.2*\x*(1-\x/6)^2}) -- cycle;
    \end{scope}
    \begin{scope}
      \clip (\c,-0.1) rectangle (6.2,3.2);
      \fill[rulegrey!20] (0,0) -- plot[domain=0:6,samples=40] (\x,{3.2*\x*(1-\x/6)^2}) -- cycle;
    \end{scope}
    \draw[ink,thick] plot[domain=0:6,samples=40] (\x,{3.2*\x*(1-\x/6)^2});
    \draw[-{Stealth}] (-0.1,0) -- (6.3,0);
    \draw[dashed,ink] (\c,0) -- (\c,3.0);
    \node[above] at (\c,3.0) {the incumbent's value};
    \node[anchor=north east] at (\c,-0.1) {to the target};
    \node[anchor=north west] at (\c,-0.1) {stays with the challenger};
    \node at (3,-1.0) {\lab};
  \end{scope}
}
\end{tikzpicture}
\caption{The picture of slices, for one incumbent value. The horizontal axis shows levels of the challenger's value, from the reserve upward. The height of the curve is the gap between the classes at that level. The shape of the curve is only a sketch. The blue area goes to the target. The grey area stays with the challenger. The total area is the same in both panels.}
\label{fig:story-slices}
\end{figure}

A stronger incumbent has a high value more often. On average its value cuts the picture further to the right. More slices go to the target, and fewer stay with the challenger. The challenger's profit falls, and the spread rises. This is Proposition~1.

The two shares always add up to the same total, whatever the incumbent. The spread plus the profit gap between the classes equals the gap in value between the classes. A stronger incumbent moves value from the challenger to the target and leaves the gap unchanged.

Test the picture at two extremes. Suppose the incumbent's value lies below every level at which the classes differ. No slice reaches the target, and the spread is zero. Suppose the incumbent's value lies above every such level. Every slice reaches the target, the challenger never wins, and the spread equals the whole gap.

\begin{watchout}
Proposition~1 uses weak words. ``Never raises'' and ``never narrows'' allow no change at all. A strict change needs a condition on the slices, which Chapter~\ref{ch:auction} states. Part~(b) also needs the class ranking. Without it, a reserve that blocks the low class more often than the high class can make the spread narrower.
\end{watchout}

The two returns pull apart. The challenger's return to preparation falls, and the investor's return to knowledge of the class rises.

{\small\color{notegrey}Source: Section 3.1 of the paper and Proposition~1.}

\section{The main result as a chain}\label{sec:story-chain}

Proposition~2 is the main result. It compares three economies: a weak incumbent, a strong incumbent that strengthens the weak one, and a very strong incumbent that strengthens the strong one. The classes, the reserve, the prior, the order slot and the costs are the same in all three. Each economy has a unique equilibrium outcome under all four tie rules. A tie rule says what an indifferent investor or an indifferent high-cost challenger does.

\begin{watchout}
The three economies are separate, and the equilibrium is solved again in each one. The result compares targets that differ only in the incumbent. It does not follow one target as its incumbent grows.
\end{watchout}

Table~\ref{tab:story-three} gives the outcome of each economy in words.

\begin{table}[t]
\centering\small
\caption{The three economies of Proposition~2, in words.}\label{tab:story-three}
\begin{tabular}{@{}>{\raggedright\arraybackslash}p{2.3cm}>{\raggedright\arraybackslash}p{3.8cm}>{\raggedright\arraybackslash}p{4.2cm}>{\raggedright\arraybackslash}p{3.7cm}@{}}
\toprule
 & Weak incumbent & Strong incumbent & Very strong incumbent\\
\midrule
The spread & Below the trading cost: a correct order earns less than it costs & Wide: a correct order covers the trading cost, whatever the investor's strategy & At least as wide as against the strong incumbent\\[2pt]
The investor & Does not trade & Class \(H\) buys and class \(L\) sells & The same orders\\[2pt]
The price & Silent: the belief stays at the prior & Informative: the belief falls after a sale and rises after a purchase & Informative: the same beliefs\\[2pt]
Gross profit & At the prior, below the high cost & At the top belief, above the high cost & At the top belief, below the high cost\\[2pt]
High-cost challenger & Stays out & Prepares after a purchase & Stays out\\[2pt]
Entry & The floor \(\rho\) & Above the floor & Back at the floor \(\rho\)\\
\bottomrule
\end{tabular}
\end{table}

Figure~\ref{fig:story-chain} draws the logic that leads from the weak to the strong incumbent. The solid chain is the information force. The arrow on the right is deterrence. The dashed box marks the link that fails against a very strong incumbent.

\begin{figure}[t]
\centering
\begin{tikzpicture}[
  box/.style={draw=rulegrey, rounded corners=2pt, fill=white, font=\small, text width=6.0cm, align=center, inner sep=4pt},
  side/.style={box, text width=5.0cm, draw=notegrey},
  turn/.style={box, text width=3.6cm, dashed, draw=notegrey, font=\footnotesize},
  flow/.style={-{Stealth[length=2.4mm]}, accent, thick},
  opp/.style={-{Stealth[length=2.4mm]}, ink, thick}]
\node[box, draw=accent] (n0) at (0,0) {A stronger incumbent};
\node[box] (n1) at (0,-1.6) {The target-payoff spread widens};
\node[box] (n2) at (0,-3.2) {Informed trade pays: a correct order covers the trading cost};
\node[box] (n3) at (0,-4.8) {The price carries the market makers' belief about the class};
\node[box] (n4) at (0,-6.4) {A purchase lifts the belief past the preparation threshold};
\node[box] (n5) at (0,-8.6) {A high-cost challenger prepares after a purchase};
\node[box, draw=accent] (n6) at (0,-10.4) {Entry rises above the floor \(\rho\). High-class ownership rises if class \(H\) outbids the incumbent};
\node[side] (r1) at (9.3,-1.6) {The challenger's gross profit falls at every belief};
\node[side] (r2) at (9.3,-3.2) {Deterrence: \mbox{preparation} pays less};
\node[turn] (t) at (6.2,-7.5) {Against a very strong incumbent, even the top belief leaves profit below the high cost};
\foreach \a/\b in {n0/n1, n1/n2, n2/n3, n3/n4, n4/n5, n5/n6} {\draw[flow] (\a) -- (\b);}
\draw[opp] (n0.east) -| (r1.north);
\draw[opp] (r1) -- (r2);
\draw[opp] (r2.south) |- (n6.east) node[pos=0.3, right=2pt, font=\footnotesize, text=notegrey, text width=2.4cm, align=left] {pushes entry down};
\draw[dashed, notegrey, -{Stealth}] (t.west) -- (0,-7.5) node[pos=0.55, above, font=\footnotesize, text=notegrey] {the link fails};
\end{tikzpicture}
\caption{The logic of Proposition~2. The chain on the left is the information force. The arrow on the right is deterrence, which lowers entry. The dashed box marks the link that fails against a very strong incumbent. Both changes of Proposition~1 are strict in the case shown.}
\label{fig:story-chain}
\end{figure}

The first link is that a wider spread makes informed trade pay. A correct order, a purchase in class \(H\) or a sale in class \(L\), earns at most the spread. Against the weak incumbent the spread is below the trading cost, so no correct order pays, and the investor stays out. Against the strong incumbent a correct order earns at least a fixed share of the spread, and this share exceeds the trading cost. So class \(H\) buys and class \(L\) sells. Section~\ref{sec:story-noise} explains the fixed share.

The second link is that informed trade makes the price informative. A purchase is likelier in class \(H\) than in class \(L\), so a purchase makes class \(H\) more likely. After a purchase the market makers raise their belief that the class is \(H\), and after a sale they lower it. The price follows the belief. The challenger sees the price, and so it learns the belief.

The third link is that a purchase lifts the belief high enough to cover the high cost. After a purchase the belief equals the top posterior, the highest belief that any order can support. Against the strong incumbent, gross profit at this belief is above the high cost. The paper says that a price recruits a high-cost challenger when this challenger prepares at that price. A purchase recruits, and a high-cost challenger prepares after it. A zero-cost challenger prepares after every price.

The fourth link is that entry rises. Entry is the floor \(\rho\) after a sale or after no order. It is one after a purchase. A purchase has positive probability, so entry is higher than the floor. The weak incumbent leaves entry at the floor.

Market makers know that a purchase brings the challenger in, and they put the gain from entry into the price. The investor cannot earn this gain. It earns only on what the market does not know, which is the class. So the chain holds when market makers anticipate the entry that each price induces.

High-class ownership needs one more condition. It is the prior probability of class \(H\), times entry in class \(H\), times the chance that a prepared class-\(H\) challenger beats the incumbent. Entry in class \(H\) rises, but a stronger incumbent can lower the chance to win. The clean rise needs class \(H\) to outbid the strong incumbent for sure. Then class \(H\) also outbids the weak incumbent, and ownership rises. The strong incumbent then takes wins only from the low class.

Entry rises first and falls later as the incumbent strengthens. Against a very strong incumbent the spread is at least as wide, so the investor still trades and the price still carries the same beliefs. But gross profit has fallen at every belief. Even the top posterior leaves gross profit below the high cost, and no price supports a higher belief. So no price recruits a high-cost challenger. Only the zero-cost challenger prepares, and entry returns to the floor.

\begin{checkyourself}
Against the very strong incumbent the spread is at least as wide as against the strong incumbent, and the investor still trades. Explain in words why entry returns to the floor.\label{q:story-turn}
\end{checkyourself}

The proposition holds under three conditions on the economies. Each condition compares one number with another.
\begin{itemize}
\item \emph{The high-cost window.} Against the weak incumbent, gross profit at the prior is below the high cost. Against the strong incumbent, gross profit at the top belief is above it.
\item \emph{The trading-cost window.} The trading cost is above the spread against the weak incumbent. It is also below the guaranteed minimum that a correct order earns against the strong incumbent (Section~\ref{sec:story-noise}).
\item \emph{The exclusion condition.} Against the very strong incumbent, gross profit at the top belief is below the high cost.
\end{itemize}
The window for the trading cost is empty unless competition widens the spread by a large enough factor. A prior far from one half makes the window narrower. The three conditions hold together on a nonempty open set of parameters. Chapter~\ref{ch:example} shows this in the example. Chapter~\ref{ch:ccc} proves the proposition.

{\small\color{notegrey}Source: Section 4.1 of the paper, Proposition~2, and the Introduction.}

\section{Why the price never reveals everything}\label{sec:story-noise}

A liquidity trader can place any order. Take a purchase. A liquidity trader buys with some probability, whatever the class. So a purchase can happen in class \(L\) as well as in class \(H\). The same holds for a sale and for no order. Every order has positive probability in both classes, and the market makers cannot tell who placed it.

An order is likelier in the class that the investor's strategy favours, but only by a bounded factor. The investor adds at most its own probability of presence to the liquidity probability. So, under every strategy of the investor, the belief after any order stays between a bottom posterior and a top posterior. For the order of no trade, the model keeps the chance that a liquidity trader buys at most one third. The challenger sees a price, and its belief is an average of the posteriors of the orders that share this price. It obeys the same bounds, so no price supports a belief above the top posterior.

\begin{analogy}[a busy market stall]
Picture a busy market stall. Many people buy for their own reasons. A friend of the owner hears good news, and she buys when the news is good. A passer-by sees a purchase and cannot tell her from any other buyer. A purchase is still a little likelier when the news is good, so it is weak evidence.

The picture breaks in three places. The informed investor buys only when the gain covers its cost. The market makers set the price by Bayes' rule, and a passer-by does not. The challenger does not see the purchase, only the price.
\end{analogy}

The doubt that remains after an order keeps the investor's edge alive. The market makers put the expected gain from entry into the price. What is left for the investor is the entry that the order brings, times the spread, times the doubt that remains about the class.

Entry is never below the floor, because a zero-cost challenger prepares after every price. The doubt never falls below a fixed level, because the posterior stays inside the bounds. So a correct order earns at least a fixed share of the spread. The paper calls this the trading bound. A wrong-signed order earns at most minus the trading cost, so it never pays.

The challenger sees the price and not the order. The price still tells it the market makers' belief when the spread is positive and the entry floor is positive. Suppose two orders have the same price. Entry after them is the same, because entry depends on the order only through the price, and it is at least the floor. At the same entry, a higher belief gives a higher price when the spread is positive, so two orders with one price have one belief. The price reveals the market makers' belief, and it never reveals the class.

\begin{checkyourself}
Suppose that liquidity traders never buy. Who could place a purchase, and what would the market makers believe after it? What would a purchase earn the investor? What does this show about the role of the liquidity order in the same direction?\label{q:story-doubt}
\end{checkyourself}

{\small\color{notegrey}Source: Sections 2.3, 3.2 and 3.3 of the paper.}

\section{Two controls}\label{sec:story-controls}

A sceptic can say that something other than information explains the rise in entry. The paper runs two controls to test this.

The first control holds the information in the price fixed while the incumbent strengthens. A stronger incumbent can then only lower entry. It lowers gross profit at every belief, so it can only remove a reason to prepare. Entry falls by exactly the probability that the preparation cost lies between gross profit against the two incumbents. This holds for every experiment of this type and every cost law that is independent of the rest.

In the order slot, the control fixes the investor's orders at their informative level. It needs a positive spread in each economy that it compares, so that competitive prices reveal the same posterior after each order. The control imposes the orders and does not solve the equilibrium again. In the example, entry with the orders held fixed stays flat from the weak to the strong incumbent. It falls only above a strength that the paper calls the entry ceiling. The control therefore removes the rise.

The second control hides the price from the challenger. The challenger then decides on its prior belief. In the example the prior lies below the preparation threshold at every strength of the path. So a high-cost challenger never prepares, and entry stays at the floor. The investor still trades against the strong incumbent. Trade alone does not raise entry.

\begin{intuition}[where the reversal comes from]
With the information fixed, entry cannot rise. With the information present and hidden from the challenger, entry stays at the floor in the example. Entry rises only when trading creates information and the challenger reads it from the price. A stronger incumbent changes the information that trading creates. The entry response to this information exceeds the direct effect of deterrence.
\end{intuition}

\begin{checkyourself}
A researcher finds that, in an economy where the challenger cannot see the price, entry rose when the incumbent became stronger. Which part of the story does this contradict, and why?\label{q:story-hidden}
\end{checkyourself}

{\small\color{notegrey}Source: Section 4.2 of the paper and the Introduction.}

\section{Quiet markets, research, and robustness}\label{sec:story-markets}

Two numbers sum up the auction for the rest of the story. Suppose that the spread is positive and class \(H\) earns more than class \(L\). Then the investor's orders, the market makers' beliefs and the challenger's choices depend on the auction only through the spread and the preparation threshold.

A stronger incumbent never narrows the spread. It never lowers the threshold while a challenger known to be in class \(L\) earns at most the high cost. So the two numbers pull in opposite directions. A wider spread helps recruitment, and a higher threshold makes recruitment harder.

Proposition~3 looks at the strengths between the weak and the strong incumbent. Against one and the same incumbent, two equilibria can exist. In a quiet market the investor does not trade, the price is silent, and entry stays at the floor. In an informative market the investor trades, and a purchase recruits a high-cost challenger. Each state confirms itself. A silent price leaves an order little to earn, and a purchase that recruits earns more.

Proposition~3 uses the default tie rules: an indifferent investor may mix between its best orders, and an indifferent high-cost challenger prepares. The two markets exist together exactly when all of these hold.
\begin{itemize}
\item The preparation threshold lies above the prior and within reach of a purchase.
\item The spread is wide enough for a purchase that just reaches the threshold to cover the trading cost.
\item The spread is narrow enough that no order against a silent price covers the trading cost.
\end{itemize}
In the informative market, entry rises with the spread while class \(H\) mixes between a purchase and no order. It stays constant once class \(H\) buys for sure. As the incumbent strengthens along a path, the economy can enter this region, leave it, and enter it again. Chapter~\ref{ch:coexistence} proves the proposition.

Proposition~4 lets the investor pay for its information. It chooses the probability that research succeeds, up to a cap, and each extra unit of research costs more than the one before. A success reveals the class and fills the slot with an order, and after a failure a liquidity trader fills it. The proposition fixes the prior at \(1/2\) and uses the default tie rules. Consider the equilibrium in which a purchase after a success recruits a high-cost challenger. In it, research and entry rise with the spread until research reaches its cap.

A wider spread raises the return from a purchase, so the investor does more research. A purchase is then more often informed, and the belief after it rises. The low class sells on bad news, and this does not stop the rise. When class \(H\) outbids the incumbent, high-class ownership rises with research. Take a threshold above the prior. The equilibrium exists exactly under two conditions.
\begin{itemize}
\item The spread exceeds twice the trading cost.
\item The belief after a purchase reaches the threshold, at the research that the investor chooses.
\end{itemize}
Chapter~\ref{ch:research} proves the proposition.

The rise in entry survives changes in the model. It survives other laws for the preparation cost, under conditions on the costs. For every prior and every pair of ranked classes that differ above the reserve, some uniform incumbents and costs give the reversal. With more than two classes, it holds under a separation condition. At the prior \(1/2\) it survives private signals that the challenger and the investor hold about different facts. Chapter~\ref{ch:robustness} states the conditions.

{\small\color{notegrey}Source: Sections 3.4, 4.3, 5 and 6 of the paper, Propositions~3 and~4.}

\section{What the paper offers for data and for the seller}\label{sec:story-uses}

A researcher can test the story with only prices and order counts, and Chapter~\ref{ch:empirical} derives the tests. The paper dates the decision interval from the public sale news to the challenger's decision. Take a market in which the investor always buys after good news and always sells after bad news, with a prior below the preparation threshold. Weight each price move by how often its order occurs. A purchase raises the price by more than a sale lowers it exactly when a purchase recruits the challenger. The difference comes from the entry part of the price gap between a purchase and a sale.

Some predictions use the class and not entry. A prepared challenger wins more often after a purchase than after a sale, because a purchase makes class \(H\) more likely. The paper states a condition for this result. It holds in any informative market, so it tests the information in the price and says nothing about recruitment. The predictions that involve the entry part of the price gap separate the mechanism from simple deterrence. The researcher must compare targets or public states that share the same liquidity and the same informed presence, because each strength is a separate economy.

The seller sets the sale terms, so the seller's choices matter. The paper does not solve the seller's design problem. It names this problem as the next step, with the whole stage of price and trade kept. It gives three results about the terms (Chapter~\ref{ch:welfare}).
\begin{itemize}
\item \emph{Access to prices.} Under the hypotheses of Proposition~2, access to prices raises expected target proceeds and acquisition surplus net of preparation costs against the strong incumbent. Shareholders gain because a holder through the sale receives the sale proceeds and not the interim price.
\item \emph{The payment rule.} Take a simple case with one value in each class, a zero reserve, and the incumbent's top value between the two class values. In a family of bargaining rules, the weight that the rule puts on the runner-up's value decides whether competition widens or narrows the spread.
\item \emph{The reserve.} A reserve that excludes the low class can make the price informative against both incumbents. If informed trade is then full and only a purchase recruits at both strengths, entry is flat in incumbent strength.
\end{itemize}

{\small\color{notegrey}Source: Sections 7, 8 and 9 of the paper.}

\section{A map of the guide}\label{sec:story-map}

Table~\ref{tab:story-map} names each chapter and the question that it answers.

\begin{table}[t]
\centering\small
\caption{The chapters of the guide and their questions.}\label{tab:story-map}
\begin{tabular}{@{}r>{\raggedright\arraybackslash}p{4.0cm}>{\raggedright\arraybackslash}p{9.8cm}@{}}
\toprule
 & Title & Question\\
\midrule
\ref{ch:story} & The story in words & What is the whole argument, before any mathematics?\\[2pt]
\ref{ch:toolkit} & The probability toolkit & Which tools from probability and game theory does the paper use, and how does each work?\\[2pt]
\ref{ch:model} & The model & What are the rules of the game, who knows what at each date, and what must an equilibrium satisfy?\\[2pt]
\ref{ch:auction} & The auction: who gets what & When the incumbent is stronger, how does the sale divide the challenger's value between the challenger and the target?\\[2pt]
\ref{ch:beliefs} & What the price reveals & What can the challenger learn from the one price it sees, and when does that belief make a high-cost challenger prepare?\\[2pt]
\ref{ch:trading} & The informed investor & What does an informed order earn, and which features of the auction does the answer use?\\[2pt]
\ref{ch:ccc} & Competition creates competition & Why can a stronger incumbent raise entry when it never raises what preparation earns at any belief?\\[2pt]
\ref{ch:example} & The example and the two controls & What does the example look like, and what do the two controls say about the source of the rise in entry?\\[2pt]
\ref{ch:coexistence} & Quiet and informative markets & Can the price stay silent in one equilibrium and recruit in another against the same incumbent, and where does this happen?\\[2pt]
\ref{ch:research} & Competition creates research & When the investor must pay for its information, does a stronger incumbent still bring informed trade and more entry?\\[2pt]
\ref{ch:robustness} & What survives when the model changes & Does the rise in entry survive other cost laws, other classes and priors, and private signals?\\[2pt]
\ref{ch:welfare} & Welfare and the terms of sale & Who gains when the challenger can see the price, and how do the payment rule and the reserve change what the price does?\\[2pt]
\ref{ch:empirical} & Taking the model to data & Which comparisons of prices, order counts and outcomes can a researcher compute, and which sign does the model predict for each?\\[2pt]
\ref{app:numbers} & The example in numbers & What are the values behind the example and its three figures?\\
\bottomrule
\end{tabular}
\end{table}

Read the guide twice. The first pass gives the mechanism. The second pass gives the proofs.

The first pass uses no algebra.
\begin{enumerate}[nosep]
\item Read this chapter.
\item In each of Chapters~\ref{ch:model} to~\ref{ch:empirical}, read the box ``In this chapter'', then the ``In words'' and ``Key idea'' boxes, then ``What to remember''. Skip the derivation boxes and the proofs.
\item Answer each ``Check yourself'' question in words, and compare your answer with the answer at the end of the chapter.
\end{enumerate}

The second pass follows the derivations.
\begin{enumerate}[nosep]
\item Skim Chapter~\ref{ch:toolkit}. Return to a section of it when a later chapter cites that section.
\item Read Chapters~\ref{ch:model} to~\ref{ch:ccc} in order. Each chapter uses the one before it.
\item Read Chapter~\ref{ch:example} next. It gives the example and the two controls of Section~\ref{sec:story-controls}.
\item Read Chapter~\ref{ch:coexistence} before Chapter~\ref{ch:research}, because Chapter~\ref{ch:research} cites it often.
\item Read Chapters~\ref{ch:robustness}, \ref{ch:welfare} and~\ref{ch:empirical} in any order. Each builds on Chapter~\ref{ch:ccc}.
\item End with Chapter~\end{enumerate}

\begin{chaptersummary}
\begin{itemize}
\item A stronger incumbent never raises the challenger's gross profit at any belief (Proposition~1(a)). This is the force of deterrence.
\item The same incumbent never narrows the target-payoff spread, which is what the informed investor's knowledge of the class is worth.
\item A wider spread makes informed trade pay, informed trade makes the price informative, and a purchase can lift the belief past the preparation threshold. Then a high-cost challenger prepares, and entry rises.
\item Against a very strong incumbent, even the top belief leaves gross profit below the high cost, so entry returns to the floor.
\item A liquidity order in the same direction leaves doubt after every order, and the doubt keeps the investor's edge alive.
\item The two controls show that the effect runs through the information that trading creates.
\item With the information fixed, entry cannot rise, and with the price hidden, entry stays at the floor in the example.
\end{itemize}
\end{chaptersummary}

\answersection

\answer{\ref{q:story-turn}}{The spread is at least as wide against the very strong incumbent, so the investor still trades. The orders and the beliefs are those of the strong incumbent. The top posterior is the highest belief that any price can support. A stronger incumbent never raises gross profit at any belief, so gross profit at the top posterior is lower than against the strong incumbent.\par In the third economy it is below the high cost. Then no price recruits a high-cost challenger. Only the zero-cost challenger prepares, and entry equals the floor. Information is still present, but the fall in profit outweighs it.}

\answer{\ref{q:story-doubt}}{Only the informed investor could place a purchase. So a purchase would show that the class is \(H\), and the market makers' belief after it would equal one. The price after a purchase would then include the whole gain that a class-\(H\) challenger brings.\par The investor earns the entry, times the spread, times the doubt that remains. The doubt would be zero, so a purchase would earn nothing, and the investor would still pay the trading cost. A purchase that shows the class therefore cannot be an informed order that pays. The liquidity order in the same direction leaves doubt after every order, and the doubt keeps the edge alive.}

\answer{\ref{q:story-hidden}}{It contradicts the claim that the effect runs through the information in the price. A challenger that cannot see the price decides on its prior. Its belief is then the same at every strength of the incumbent. A stronger incumbent never raises gross profit at a fixed belief (Proposition~1(a)). A high-cost challenger therefore prepares at most as often against the stronger incumbent, and entry cannot rise. This is the fixed-information control with a trivial experiment.}
